% from ega4-1-fr.tex
\oldpage[IV-1]{232}

\begin{proof}
On peut se borner au cas où $Z=\operatorname{Spec}(A)$, $Y=\operatorname{Spec}(B)$, $B$ étant une
$A$-algèbre de type fini~; le morphisme $\Delta_g$ correspond alors à
l'\emph{homomorphisme d'augmentation} $\delta:B\otimes_AB\to B$ tel que
$\delta(x\otimes y)=xy$. Compte tenu de la définition (1.4.2), il suffit de
montrer que le noyau $\mathfrak J$ de $\delta$ est un idéal \emph{de type
fini}, ce qui résulte de (0, 20.4.4).
\end{proof}

\begin{proposition}[1.4.4]
\label{IV.1.4.4-fr}
Soient $A$ un anneau, $B$ une $A$-algèbre, $B'$ une algèbre de polynômes
$A[T_1,\ldots,T_n]$, $u:B'\to B$ un homomorphisme surjectif de
$A$-algèbres. Pour que $B$ soit une $A$-algèbre de présentation finie, il faut
et il suffit que le noyau $\mathfrak J$ de $u$ soit un idéal de type fini de
$B'$.
\end{proposition}

\begin{proof}
La condition est suffisante par définition (1.4.1). Pour voir qu'elle est
nécessaire, remarquons que le morphisme $g:\operatorname{Spec}(B')\to\operatorname{Spec}(A)$ est
localement de type fini~; si $B$ est une $A$-algèbre de présentation finie,
le morphisme $f:\operatorname{Spec}(B)\to\operatorname{Spec}(B')$ correspondant à $u$ est tel que
$g\circ f:\operatorname{Spec}(B)\to\operatorname{Spec}(A)$ soit localement de présentation finie, donc il
résulte de (1.4.3, (v)) que $f$ est aussi localement de présentation finie.
Or, pour démontrer que l'idéal $\mathfrak J$ est de type fini, il suffit de
prouver que $\widetilde{\mathfrak J}$ est un $\widetilde{B'}$-Module de type
fini (II, 6.1.4.1). Revenant à la définition (1.4.2), on est finalement ramené
à prouver le
\end{proof}

\begin{lemma}[1.4.4.1]
\label{IV.1.4.4.1-fr}
Soient $\mathfrak a$ un idéal d'un anneau $C$, $s$ un élément de $C$, tel que
$C_s/\mathfrak a_s$ soit une $C$-algèbre de présentation finie~; alors
$\mathfrak a_s$ est un idéal de type fini dans l'anneau $C_s$.
\end{lemma}

\begin{proof}
Par hypothèse, il existe une $C$-algèbre de polynômes
$C'=C[T_1,\ldots,T_n]$ et un $C$-homomorphisme surjectif
$v:C'\to C_s/\mathfrak a_s$ dont le noyau $\mathfrak b$ est de type fini.
Soit $p:C_s\to C_s/\mathfrak a_s$ l'homomorphisme canonique~; pour tout $i$,
il y a un $t_i\in C$ tel que $p(t_i/s^k)=v(T_i)$ (on peut supposer que
l'exposant $k$ est le même pour tous les indices $i$). Considérons alors la
$C_s$-algèbre $D=C_s[T_1,\ldots,T_n]$, et l'homomorphisme de $C_s$-algèbres
$w:D\to C_s$ tel que $w(T_i)=t_i/s^k$~; $w$ est évidemment surjectif, et il
en est par suite de même de $v'=p\circ w:D\to C_s/\mathfrak a_s$. D'autre
part, tout polynôme de $D$ peut s'écrire $P/s^m$, où
$P\in C'=C[T_1,\ldots,T_n]$, et l'on a
$v'(P/s^m)=(1/s^m)v(P)$~; comme $1/s$ est inversible dans $C_s$, la relation
$v'(P/s^m)=0$ dans $C_s/\mathfrak a_s$ équivaut à $v(P)=0$, et par suite le
noyau $\mathfrak b'$ de $v'$ est engendré par l'image canonique de
$\mathfrak b$ dans $D$, et \emph{a fortiori} est de type fini dans l'anneau
$D$. Mais
$\mathfrak b'=v'^{-1}(0)=w^{-1}(p^{-1}(0))=w^{-1}(\mathfrak a_s)$, et comme
$w$ est surjectif, $\mathfrak a_s=w(\mathfrak b')$, ce qui prouve que
$\mathfrak a_s$ est un idéal de type fini de $C_s$.
\end{proof}

\begin{corollary}[1.4.5]
\label{IV.1.4.5-fr}
Soient $X$, $Y$ deux préschémas, $j:X\to Y$ une immersion, $U$ un ouvert de
$Y$ tel que $j(X)$ soit fermé dans $U$, $\mathcal J$ l'idéal quasi-cohérent de
$\mathcal O_U$ qui définit le sous-préschéma fermé de $Y$ associé à $j$. Pour
que $j$ soit localement de présentation finie, il faut et il suffit que
$\mathcal J$ soit un $\mathcal O_U$-Module de type fini.
\end{corollary}

\par\medskip\noindent
\begin{proposition}[1.4.6]
\label{IV.1.4.6-fr}
Soit $\varphi:A\to B$ un homomorphisme d'anneaux. Pour que le morphisme
correspondant $f:\operatorname{Spec}(B)\to\operatorname{Spec}(A)$ soit localement de présentation finie,
il faut et il suffit que $B$ soit une $A$-algèbre de présentation finie.
\end{proposition}

\begin{proof}
La condition étant trivialement suffisante, il reste à prouver qu'elle est
nécessaire. Si $f$ est localement de présentation finie, il résulte déjà de
(1.3.6) que $B$ est une $A$-algèbre de type fini, donc il existe un
homomorphisme surjectif $u:B'=A[T_1,\ldots,T_n]\to B$ de $A$-algèbres~; il
résulte de (1.4.3, (v)) que l'immersion
$j:\operatorname{Spec}(B)\to\operatorname{Spec}(B')$ est localement

